%------------------------------------------------

\section{Enfoque Práctico}

%------------------------------------------------

\begin{frame}{¿Cómo utiliza el compilador VFP y NEON?}
 \begin{textblock*}{\textwidth}(0.4cm,1.2cm)
  Cuando el programador escribe código C/C++ convencional, el compilador puede generar automáticamente instrucciones:
  \vspace{0.3cm}
  \begin{itemize}[<+(1)->]\justifying \parskip 0.2cm
   \item ARM enteras.
   \item VFP para operaciones en punto flotante.
   \item NEON para operaciones SIMD.
  \end{itemize}
  \vspace{0.4cm}
  \only<5>{
   \begin{tcolorbox}[enhanced jigsaw,rounded corners, drop fuzzy shadow=ShadowColor]
    \justifying Cuando el objetivo es aprovechar el hardware disponible prescindiendo de escribir código ensamblador, hay alternativas.
   \end{tcolorbox}
   }
 \end{textblock*}
\end{frame}

%------------------------------------------------

\begin{frame}{Tres niveles de programación}
 \begin{textblock*}{\textwidth}(0.4cm,1.2cm)
  Existen tres formas habituales de utilizar NEON.
  \vspace{0.5cm}
  \begin{enumerate}[<+(1)->]
   \item Código C/C++ convencional.
   \item Intrinsics NEON.
   \item Ensamblador.
  \end{enumerate}
  \vspace{0.5cm}
  \only<5>{
   \begin{tcolorbox}[enhanced jigsaw,rounded corners, drop fuzzy shadow=ShadowColor]
    \justifying A mayor nivel de abstracción disminuye el esfuerzo de programación, mientras que el ensamblador ofrece el máximo control sobre las instrucciones generadas.
   \end{tcolorbox}
   }
 \end{textblock*}
\end{frame}

%------------------------------------------------

\begin{frame}[fragile]{Auto-vectorización}
 \begin{textblock*}{\textwidth}(0.4cm,1.2cm)

El compilador puede detectar operaciones repetitivas sobre arreglos y generar automáticamente instrucciones SIMD.
\vspace{0.2cm}

\begin{lstlisting}[language=C,basicstyle=\scriptsize,numbers=none, escapechar=\@]
for (i=0; i<n; i++)
    c[i] = a[i] + b[i];
\end{lstlisting}

\vspace{0.2cm}

Con las opciones de optimización adecuadas, GCC o Clang pueden reemplazar el
bucle por instrucciones NEON.
 \end{textblock*}
\end{frame}

%------------------------------------------------

\begin{frame}[fragile]{Intrinsics NEON}
 \begin{textblock*}{\textwidth}(0.4cm,1.2cm)
  \begin{tcolorbox}[enhanced jigsaw,rounded corners, drop fuzzy shadow=ShadowColor]
   \justifying Los \emph{intrinsics} son funciones de C que representan casi directamente una instrucción NEON.
  \end{tcolorbox}

  \vspace{0.3cm}

  \begin{lstlisting}[language=C,basicstyle=\scriptsize,numbers=none, escapechar=\@]
uint8x16_t va, vb, vc;

vc = vaddq_u8(va, vb);
  \end{lstlisting}

  \vspace{0.3cm}
  El compilador traduce esta llamada a la instrucción NEON correspondiente, sin necesidad de recurrir a un archivo fuente adicional con código assembler.
 \end{textblock*}
\end{frame}

%------------------------------------------------

\begin{frame}{¿Cuándo utilizar intrinsics?}
 \begin{textblock*}{\textwidth}(0.4cm,1.2cm)\justifying
  Los intrinsics permiten:
  \vspace{0.5cm}
  \begin{itemize}[<+(1)->]\parskip 0.2 cm
   \item Obtener un mayor control sobre el código generado.
   \item Acceder a instrucciones específicas de NEON.
   \item Mantener la portabilidad entre compiladores compatibles.
  \end{itemize}
  \vspace{0.4cm}
  \only<5>{En la mayoría de las aplicaciones representan el mejor compromiso entre rendimiento y facilidad de mantenimiento.}
 \end{textblock*}
\end{frame}

%------------------------------------------------

\begin{frame}{¿Y el ensamblador?}
 \begin{textblock*}{\textwidth}(0.4cm,1.2cm) \justifying
  El lenguaje ensamblador sigue siendo útil cuando:
  \vspace{0.4cm}
  \begin{itemize}[<+(1)->]\parskip 0.2 cm
   \item se requiere el máximo rendimiento;
   \item se desea controlar exactamente las instrucciones emitidas;
   \item es necesario aprovechar características muy específicas del hardware.
  \end{itemize}
  \vspace{0.4cm}
  \only<5>{Sin embargo, con el transcurrir del tiempo los compiladores perfeccionaron el uso de \emph{intrinsics}, con lo cual la utilización de ensamblador es menos frecuente.}
 \end{textblock*}
\end{frame}
